\section{Análisis de resultados}
\label{sec:resultados}

\subsection{Las dos estrategias son igual de correctas}

En las cuatro carreras de escritura, bajo cien usuarios durante sesenta
segundos y con cuatro procesos compitiendo, la verificación contra PostgreSQL
dio siempre lo mismo: suma global del libro mayor igual a cero, cero cuentas
descuadradas y cero saldos negativos. Ese es el resultado que el proyecto
quería demostrar, y es independiente de la estrategia empleada.

Dicho de otro modo: las garantías no dependen de haber elegido bien entre
bloqueo pesimista y \codigo{SERIALIZABLE}. Dependen de que las invariantes
vivan en el motor y de que la transferencia sea una sola transacción.

\subsection{Se diferencian en cuánto trabajo terminan}

Bajo contención máxima, la estrategia pesimista confirma 148,3 transferencias
por segundo y \codigo{SERIALIZABLE} solo 10,4. Catorce veces menos, aunque su
tasa de peticiones por segundo (63,1 frente a 147,0) sugiera una diferencia
mucho menor. La explicación está en la columna de fallos: el 83,72\,\% de las
peticiones termina en un 503 porque se agotan los tres reintentos ante
\sqlstate{40001}.

\subsection{La conclusión depende de dónde se mire}

Este es el hallazgo metodológico de la fase. Si el análisis se hubiera quedado
en las peticiones por segundo que reporta el generador de carga, la conclusión
habría sido \enquote{\codigo{SERIALIZABLE} rinde un 57\,\% menos}. Contando
las filas realmente escritas en PostgreSQL, la conclusión es que hace un
93\,\% menos de trabajo útil.

Son dos afirmaciones muy distintas sobre el mismo sistema, y solo una es
cierta. Medir la respuesta de la API y verificar que las transacciones se
registraron no son la misma cosa, y confundirlas es fácil.

\subsection{Comparación con la literatura}

Ports y Grittner reportan para su implementación un coste inferior al 7\,\%
respecto al aislamiento por instantáneas~\cite{ports2012}. Las mediciones de
PayCore bajo contención máxima arrojan una diferencia mucho mayor. No hay
contradicción, y la razón importa:

\begin{itemize}[leftmargin=1.2em]
  \item Sus cargas de evaluación reparten el trabajo entre muchas filas. El
    escenario \codigo{caliente} de PayCore concentra \emph{todas} las
    escrituras en una sola fila, que es el peor caso concebible para una
    estrategia optimista.
  \item Su cifra mide el coste de la \emph{técnica} cuando los conflictos son
    raros. La de PayCore mide el efecto de los \emph{conflictos} cuando son
    constantes.
  \item En el escenario de contención baja, con cien usuarios sobre cincuenta
    cuentas, la diferencia se reduce pero no desaparece: sigue habiendo un
    27,35\,\% de fallos. Haría falta una dispersión mucho mayor para que
    \codigo{SERIALIZABLE} compitiera, y ese no es el perfil de carga de una
    billetera.
\end{itemize}

\subsection{Verificación de la hipótesis previa}

El registro de decisión sobre concurrencia declaró, \emph{antes} de medir, que
bajo contención máxima se esperaba que la estrategia optimista degenerase en
abortos y reintentos. La medición lo confirma. La hipótesis estaba escrita con
fecha anterior a los datos, lo que permite distinguir una predicción acertada
de una racionalización posterior.

\subsection{Qué aportó cada tipo de prueba}

\begin{table}[htbp]
\centering
\caption{Defectos reales encontrados, y por qué tipo de prueba.}
\label{tab:hallazgos}
\small
\begin{tabularx}{\textwidth}{@{}lXX@{}}
\toprule
\textbf{Tipo} & \textbf{Qué encontró} & \textbf{Por qué nada más lo veía} \\
\midrule
Unitarias y de integración & Nada que las de concurrencia no vieran después &
  Un solo hilo nunca reproduce una carrera \\
Concurrencia & La actualización perdida del mapa de identidad: 88
  transferencias aprobadas de 100 &
  Las 76 pruebas existentes pasaban todas con el defecto presente \\
Carga & El agotamiento del pool de conexiones por retener la sesión fuera del
  hilo & Solo aparece con suficientes peticiones simultáneas y rápidas \\
Revisión de la propia medición & Una carrera truncada a 34 s que invalidaba
  una comparación & Los totales parecían razonables; solo las marcas de
  tiempo lo delataban \\
\bottomrule
\end{tabularx}
\end{table}

La lectura de este cuadro es la conclusión práctica del proyecto: cada tipo de
prueba encuentra una clase de defecto que las demás no pueden ver, y eso
incluye revisar críticamente las propias mediciones.
